% !TeX program = lualatex
% This document requires LuaLaTeX for compilation

\documentclass[a4paper,12pt]{article}

% Engine detection and compatibility check
\RequirePackage{iftex}
\ifLuaTeX
  % Document is being compiled with LuaTeX - proceed normally
\else
  \PackageError{main}{This document requires LuaLaTeX for compilation.
    Please use: lualatex main.tex}
    {This document uses fontspec and other LuaTeX-specific features.}
\fi

\usepackage{fontspec}
\usepackage[brazilian]{babel}

% Font configuration for LuaTeX
\setmainfont{Latin Modern Roman}
\setsansfont{Latin Modern Sans}
\setmonofont{Latin Modern Mono}

\usepackage{amsmath} % Para equações matemáticas
\usepackage{geometry}
\geometry{margin=1in}
\usepackage{hyperref}
\usepackage[sort,numbers]{natbib}
\usepackage{parskip}
\usepackage{longtable}
\usepackage{xcolor}
\usepackage{tabularx}
\setlength{\marginparwidth}{2cm}
\usepackage[colorinlistoftodos,prependcaption,textsize=tiny]{todonotes}
\usepackage{comment}
\usepackage{graphicx} % Para resizebox
\usepackage{array} % Para colunas com largura fixa
\usepackage{listings}
\usepackage{siunitx} % For \si{\micro\second}
\usepackage[acronym,nonumberlist]{glossaries}
\usepackage{placeins}

% Configuring listings for code and log display
\lstset{
  basicstyle=\ttfamily\small,
  breaklines=true,
  breakatwhitespace=true,
  frame=single,
  numbers=left,
  numberstyle=\tiny,
  keywordstyle=\color{blue},
  commentstyle=\color{gray},
  showstringspaces=false,
  tabsize=2,
  extendedchars=true
}


\title{Trabalho 2 - Algoritmos Genéticos - Problema das Oito Rainhas}
\author{André Thiago Pfleger \\  Gustavo Girotto \\ João Pedro Schmidt Cordeiro}
\date{\today}

\begin{document}

\begin{titlepage}
   \begin{center}
      \textsc{Universidade Federal de Santa Catarina} \\
      \textsc{Centro Tecnológico} \\
      \textsc{Departamento de Informática e Estatística} \\[4cm]

      \textbf{\Large{INE5430 - Inteligência Artificial}} \\[2cm]

      \textbf{\huge{Trabalho 2 - Algoritmos Genéticos - Problema das Oito Rainhas}} \\[3cm]

      André Thiago Pfleger \\
      Gustavo Girotto \\
      João Pedro Schmidt Cordeiro \\
      
      \vfill
      
      {Florianópolis} \\
      {\today}
   \end{center}
\end{titlepage}

\section{Introdução}
\label{ch:introducao}

Este trabalho apresenta o desenvolvimento de um algoritmo genético para encontrar a solução do problema da oito rainhas.

O objetivo principal é demonstrar a aplicação prática de algoritmos genéticos para encontrar a solução do problema da oito rainhas.


\section{Problema}
\label{ch:problema}

O problema das oito rainhas é um quebra-cabeça matemático e um exemplo clássico de problema de satisfação de restrições, que consiste em posicionar oito rainhas em um tabuleiro de xadrez de 8x8 casas de tal forma que nenhuma rainha possa atacar a outra. Para que isso seja possível, não pode haver duas ou mais rainhas na mesma linha, coluna ou diagonal \cite{wiki:8queens}.

O problema foi originalmente proposto em 1848 pelo enxadrista Max Bezzel. Desde então, tem sido estudado por muitos matemáticos e cientistas da computação, servindo como um problema comum para testar algoritmos de busca e otimização.

\section{Implementação}
\label{ch:implementacao}

A implementação do algoritmo genético foi feita em Python usando a biblioteca \texttt{pygad} \cite{pygad}.

\subsection{Cromossomo}
\label{sec:cromossomo}

O cromossomo é a representação do problema. Neste caso, o cromossomo é uma lista de 8 inteiros, onde cada inteiro representa a coluna da rainha na linha correspondente.

\subsection{População}
\label{sec:populacao}

A população é o conjunto de cromossomos. Neste caso, a população é um conjunto de 50 cromossomos.

\subsection{Função de Fitness}
\label{sec:funcao_fitness}

A função de fitness é a função que avalia a qualidade de um cromossomo. Neste caso, a função de fitness é a soma do número de pares de rainhas que se atacam.

\subsection{Critério para Seleção Cruzamento e Mutação}
\label{sec:critario_selecao_cruzamento_mutacao}

<Critério para seleção cruzamento e mutação(e as respectivas taxas)>

\subsection{Critério de Parada e Plots de Execução}
\label{sec:critario_parada_plots_execucao}

<Critério de parada e plots de execução, mostrando a convergência do algoritmo e melhor indivíduo da população.>

O critério de parada é a condição que determina quando o algoritmo deve parar. Neste caso, o algoritmo deve parar quando encontrar uma solução satisfatória.

\section{Conclusão}
\label{ch:conclusao}

<Considerações sobre a implementação, dificuldades encontradas e percepções pessoais.>

\bibliographystyle{IEEEtran}
\bibliography{ref}

\end{document}
